\chapter{Consciousness, and the World It Comes With}
\label{sec:consciousness}
% ===================================================================

The last four chapters built a lattice and argued that every agent
occupies a point in it.
This chapter says what occupying a point amounts to.

The claim is not that high positions are better at thinking.
It is that agents at different positions \emph{think differently and
inhabit different worlds}, and that the second half of that sentence is
the substantive one: the ontology an agent's bisimulation generates at a
point is strictly richer than the ontology generated at any point below
it.
Consciousness, in this framework, is not something an agent has more or
less of.
It is where the agent is, and where it is fixes what there is.

\section{The proposal}

\begin{definition}[Consciousness]
\label{def:consciousness}
Let $\mathfrak{W}$ be the Weihrauch lattice and let
$\mathrm{pos}(A) \in \mathfrak{W}$ be the choice strength agent $A$ can
afford to resolve at its cuts, in the sense of Chapter~\ref{ch:agy}.
The \emph{consciousness} of $A$ is $\mathrm{pos}(A)$.
An agent is \emph{conscious at} $w$ when $\mathrm{pos}(A) = w$: it can
form and evaluate the queries that $w$ resolves, and its world model
reflects the bisimulation quotient $\bisim_w$ available at $w$.
\end{definition}

An agent that resolves only what finite contention resolves --- which
earlier drafts of this book called \emph{sentient but not conscious} ---
is an agent at the bottom of the lattice.
It is not outside the structure.
It is at the floor of it.

\begin{remark}[Consciousness is a coordinate, not a threshold]
\label{rmk:coordinate}
The tower of Chapter~\ref{ch:tower} is indexed by ordinals, and
\S\ref{sec:tower-wrong} has already said, in that chapter, that the index
is wrong.
Definition~\ref{def:consciousness} takes the index from the lattice
instead, and it says something the ordinal version cannot.

Every agent occupies a position, and the position \emph{is} its
consciousness.
There is no threshold separating the conscious from the non-conscious,
and there is no transfinite ladder of grades either; there is a lattice,
and everything is somewhere in it.
The question ``is this system conscious?'' does not dissolve into ``at
what level?'' --- it dissolves into \emph{where}, and the answer is a
point that need not be comparable with the point another agent occupies.
Two agents can each resolve what the other cannot.
An ordinal index cannot say that; a lattice coordinate says it without
further apparatus.

The second thing the coordinate reading makes explicit is that an agent's
position is simultaneously the strength at which the world \emph{islands}
for it, in the sense of the graded factorization of
Chapter~\ref{ch:agy}.
The distinctions an agent can resolve and the granularity of the world it
finds itself in are not two facts that happen to correlate.
They are one fact, reported once from the inside and once from the
outside.
\end{remark}

\section{Consciousness is a resource of a place, not a property of an agent}
\label{sec:shared}

Definition~\ref{def:consciousness} attaches a lattice point to an agent,
which makes it look like a property the agent owns.
It is worth being explicit that it is not, because almost everything
interesting in this chapter follows from its not being.

\begin{proposition}[Co-located agents share what they can resolve]
\label{prop:shared}
Let $A$ and $B$ be agents with $\mathrm{pos}(A) = \mathrm{pos}(B) = w$.
Then the class of scheduling problems $A$ can resolve and the class $B$
can resolve are the same class, and the bisimulation quotient each
world model can distinguish to is the same quotient $\bisim_w$.
\end{proposition}

\begin{proof}
Both statements are functions of $w$ alone.
The first is the definition of a Weihrauch degree; the second is the
construction of $\bisim_w$ from what the degree resolves, which mentions
no agent.
\end{proof}

\begin{remark}[What is shared, and what is not]
\label{rmk:shared}
Proposition~\ref{prop:shared} is nearly trivial to prove and its content
is entirely in what it licenses one to say.

Consciousness in this framework is not a private glow that each agent
carries around.
It is an \emph{affordance of a position}: what can be resolved there,
hence what can be distinguished there, hence what exists there.
Every agent standing at $w$ draws on the same affordance, in the way that
every organism in a habitat draws on the same light.
Two agents at $w$ inhabit the same world, in the strong sense that the
inventory of things in it is the same inventory.

What is emphatically not shared is $\suff$.
Suffering, by Definition~\ref{suf:def:suffering}, is a reading an
individual takes of its own reserves against its own setpoint, and two
agents at the same lattice point can be in wholly different states of it.
That is the cleanest form of the distinction this chapter is drawing:
\emph{sentience is owned and consciousness is occupied}.
One is a fact about a learner's account; the other is a fact about a
neighborhood in a lattice which the learner happens to be standing in.
\end{remark}

\begin{remark}[Why the ecology could not have said this]
\label{rmk:ecology-couldnt}
Part~\ref{part:ecology} treats a population as a distribution of learners
over a namespace, and its notion of what several learners have in common
is a shared \emph{reservoir}.
That is a resource in the ordinary sense: rival, exhaustible, and reduced
by each learner's draw.

The resource in Proposition~\ref{prop:shared} is not like that at all.
$A$'s resolving a query does not consume $B$'s capacity to resolve it,
and a crowd at $w$ does not make $w$ poorer.
A framework whose only notion of sharing was the reservoir had no way to
express a common possession of this kind, which is one reason this
chapter is here and not there.
\end{remark}

\section{A different world, and strictly a richer one}
\label{sec:richer}

We can now argue the claim the chapter opened with.
Throughout this section, $v < w$ is the lattice order: $w$ resolves
everything $v$ resolves and something more.
The old presentation of this material stated each result for an ordinal
successor $\alpha \to \alpha+1$; every one of them is really a statement
about the order, and that is how they are stated here.

\subsection{Finer bisimulation quotient}

\begin{proposition}[Strict refinement of bisimulation]
\label{prop:strict-bisim}
Let $v < w$ in $\mathfrak{W}$ and let $\GSLT$ be Turing complete.
Then $\bisim_w \subsetneq \bisim_v$ as equivalence relations on terms:
there exist $P, Q$ with $P \bisim_v Q$ and $P \not\bisim_w Q$.
\end{proposition}

\begin{proof}[Proof sketch]
By diagonalization.
Since $v < w$ there is a scheduling problem resolved at $w$ and not at
$v$.
Take $P, Q$ whose distinguishability turns on the answer: their observable
behaviors coincide on every experiment an agent at $v$ can complete, and
differ on one which requires that resolution.
Then $P \bisim_v Q$, and an agent at $w$ separates them in one
interaction.
\end{proof}

\begin{remark}[The atoms of behavior are different at different places]
The proposition says that the \emph{atoms of observable behavior} are not
the same at $v$ and at $w$.
Terms that are experimentally indistinguishable at $v$ --- that no agent
there can tell apart, however long it experiments --- come apart at $w$.
There are more things in the world at $w$.
Not more information about the same things: more things.
\end{remark}

\subsection{Strictly more expressive hypothesis language}

\begin{proposition}[Strict expansion of the hypothesis language]
\label{prop:strict-hml}
For $v < w$, $\HML(\ctx_v) \subsetneq \HML(\ctx_w)$: there are formulae
expressible at $w$ and not at $v$.
\end{proposition}

\begin{proof}[Proof sketch]
Resolution at $w$ gives rise to minimal contexts not present at $v$ ---
those placing a term in interaction with what $w$ resolves and $v$ does
not --- and these generate modal operators absent from $\HML(\ctx_v)$.
\end{proof}

\begin{remark}
An agent at $v$ not only cannot \emph{test} certain hypotheses; it cannot
\emph{form} them.
The concepts needed to think the thought are absent from the hypothesis
language.
This is the sense in which the two agents think differently: not more
slowly and more quickly, but in different vocabularies, one of which
contains words the other has no translation for.
\end{remark}

\subsection{New conserved quantities}

\begin{proposition}[New currents higher in the lattice]
\label{prop:new-noether}
Weight and cost maps at $w$ can be sensitive to distinctions invisible at
$v$.
Symmetries of the $w$-cost map which are not symmetries of the $v$-cost
map give rise to conserved quantities at $w$ with no analogue at $v$.
\end{proposition}

\begin{remark}
In physical terms the world at $w$ has strictly more conservation laws.
What appears at $v$ as one undifferentiated interaction resolves at $w$
into a structured process with internal symmetries and corresponding
currents.
The analogy with the hierarchy of effective field theories is direct:
new conservation laws --- baryon number, lepton number, color charge ---
appear as one descends to finer scales.
\end{remark}

\subsection{Richer internal causal structure}

\begin{proposition}[Resolution of elementary interactions]
\label{prop:resolution}
A transition $P \rewrite_v Q$ that is a single step at $v$ may resolve at
$w$ into a structured history
$P = R_0 \rewrite_w \cdots \rewrite_w R_k = Q$
with internal causal structure visible only from $w$.
\end{proposition}

\begin{remark}[What existed turns out to be composite]
\label{rmk:composite}
This is the sharpest form of the richness claim, and it is worth stating
without the mathematics.
What looked like an atom of causal history at $v$ has parts at $w$.
Not merely more of them: the things that were already there turn out to
have insides.
The analogy with the discovery that atoms have internal structure, and
nuclei theirs, and protons theirs, is direct and intended.

And it is the fact this act will need in Chapter~\ref{ch:sim}.
An agent at $v$ looking at $w$'s world does not see a world with things
missing from it.
It sees a world of the right size whose contents happen to be
featureless, and it has no way, from inside, to tell that reading from
the true one.
\end{remark}

\section{Sentience and consciousness}

The distinction can now be stated exactly.

\begin{center}
\renewcommand{\arraystretch}{1.4}
\small
\begin{tabular}{lll}
\toprule
& \textbf{Sentience} & \textbf{Consciousness} \\
\midrule
What it is & A reading a learner takes & A position a learner occupies \\
Of what & Its account, against a setpoint & The choice strength it affords \\
Type & Signed vector over flavors & A point in a lattice \\
Owned or shared & Owned; private to the learner & Shared by all co-located \\
Changes & With inflow, outflow, and setpoint & By moving in the lattice \\
Analogue & Affect & The world one is in \\
Where built & Chapter~\ref{ch:suffer} & This chapter \\
\bottomrule
\end{tabular}
\end{center}

\begin{remark}[The two vary independently]
\label{rmk:independent}
An agent at the floor of the lattice can have a rich and violently
changing $\suff$.
Its inflow fails, its drift is read, its policy shifts, and all of this
is causally efficacious.
But the world it lives in is the coarse world of the floor: it lacks the
concepts to form the hypotheses available higher up, and those
distinctions are simply not in its experience.

Conversely, nothing about standing high in the lattice makes an agent
feel anything.
An agent with no setpoint at all has $\suff$ undefined and may sit
anywhere.
Nothing about being at the floor makes an agent numb; nothing about being
high makes it tender.
The two quantities are independent, and keeping them apart is the whole
reason for two words.
\end{remark}

\begin{remark}[Can a system be conscious without being sentient?]
Definition~\ref{def:consciousness} does not require a setpoint, so the
answer is formally yes.
Whether it is stable is another matter:
Conjecture~\ref{suf:conj:selected} says internalization is selected
wherever inflow is ragged, and a position high in the lattice is of no use
to a lineage that starves before it can occupy one.
Consciousness without sentience is possible in principle and, in any
environment with variance, unlikely to persist.
\end{remark}

\section{The two cycles}
\label{sec:relationship}

The interaction between the two runs in both directions, and it is the
reason they were ever confused.

\paragraph{The favorable cycle.}
Resolution at $w$ gives a finer world model.
A finer world model gives sharper anticipation of inflow, which by
Remark~\ref{ent:rmk:thermostat} is the entire point of paying for
regulation.
Sharper anticipation holds $\vect{A}$ nearer $\setpt$, so $|\suff|$ falls
and more of the account is free.
A freer account funds the assays that keep resolution at $w$ affordable
--- because by Chapter~\ref{ch:agy} occupying a position is not free
either; it is a standing charge.
And so back to the beginning.

\begin{remark}
The cycle is not guaranteed.
It is broken by environmental disruption, by exhaustion, and by
adversarial interaction with other agents, of which
Part~\ref{part:ecology} supplies a great deal.
What it describes is the stable attractor of a well-funded learner in a
benign environment, and the name for that in another vocabulary is
flourishing.
\end{remark}

\paragraph{The unfavorable cycle.}
Inflow falls.
Drift grows, $|\suff|$ grows with it, and by
Proposition~\ref{suf:prop:policy} the learner reallocates its search
toward the components in deficit --- which is correct, and which is
paid for by giving up assays about the world.
The world model coarsens.
Anticipation degrades, so drift grows further.
At some point the standing charge of the learner's position stops being
affordable, and the learner does not lose consciousness so much as
\emph{descend}: it comes to occupy a lower point, where fewer things
exist.

\begin{remark}[Descent is ontological, not merely cognitive]
\label{rmk:descent}
This is the least comfortable consequence in the chapter and it should
not be softened.
A learner that can no longer afford its position does not merely become
worse at seeing the world it was in.
By Proposition~\ref{prop:strict-bisim} it comes to be in a different
world, one with fewer things in it, and by
Proposition~\ref{prop:strict-hml} it loses the vocabulary in which the
missing things could be missed.
It does not experience the loss as loss.
There is nothing left to notice it with.

The structural parallel to what is called cognitive decline is
uncomfortably exact, and we make no clinical claim whatever.
\end{remark}

\section{Graded position}

The lattice is not discrete, and an agent's ability to resolve a
scheduling problem may be partial.

\begin{definition}[Approximate resolution]
\label{def:approx}
An agent $A$ has \emph{$\epsilon$-approximate resolution at $w$} if for
each instance of the scheduling problem characteristic of $w$ it returns
a correct answer with probability at least $1-\epsilon$, where the
probability is the one supplied by the stochastic instance of the weight
construction of Chapter~\ref{ch:weight}.
\end{definition}

\begin{remark}[A correction, and what it costs]
\label{rmk:approx-status}
Earlier presentations of this definition measured the probability as
$|\langle \mathtt{halt} \mid A \inter \lceil (M,x) \rceil \rangle|^2$ ---
a squared amplitude, on the complex instance.
That instance was proposed and withdrawn in
Chapter~\ref{sec:quantum}, and by Proposition~\ref{wt:prop:interference}
an amplitude attached to a rewrite carries no coherence in any case.
The quantity available is a \emph{rate} and not an amplitude unless the
presentation supplies a Hamiltonian, and
Definition~\ref{def:approx} is stated on that basis.

What is lost by the correction is the continuous interpolation the
amplitude version advertised.
A rate gives a probability of being right, which grades an agent's
\emph{reliability} at $w$; it does not by itself place the agent at a
point strictly between $v$ and $w$.
Whether approximate resolution defines a genuine intermediate position in
$\mathfrak{W}$ is open, and is the honest residue of a claim that used to
be made too easily.
\end{remark}

% ===================================================================
